\chapter{PENDAHULUAN}
\label{bab:pendahuluan}

Bab ini memperlihatkan susunan pendahuluan pada contoh naskah skripsi. Seluruh uraian merupakan petunjuk penulisan dan contoh format yang perlu diganti dengan materi penelitian pengguna. Naskah ini tidak mewakili proposal yang telah disetujui, penelitian yang telah dilaksanakan, ataupun hasil pengujian suatu metode. Tujuan contoh adalah memperlihatkan hubungan antara bagian naskah, penomoran, dan penyajian unsur akademik dalam satu dokumen yang dapat dikompilasi.

\section{Latar Belakang}

Pendahuluan membantu pembaca mengenali masalah, konteks, serta alasan suatu penelitian perlu dilakukan. Paragraf awal dapat memperkenalkan keadaan umum yang relevan dengan objek penelitian. Paragraf berikutnya mempersempit konteks hingga persoalan yang dapat diamati dan diteliti. Alur tersebut memungkinkan pembaca memahami hubungan antara kebutuhan penelitian dan pendekatan yang dipilih, tanpa harus menebak latar pemilihan topik.

Dalam naskah penelitian sesungguhnya, uraian latar belakang perlu ditopang sumber yang telah dibaca dan diperiksa. Pernyataan tentang keterbatasan pendekatan terdahulu harus sesuai dengan isi sumber yang dirujuk. Jika keterbatasan hanya berlaku pada objek atau kondisi tertentu, batas tersebut perlu dipertahankan dalam kalimat. Sitasi tidak menggantikan penjelasan mengenai bagaimana sumber mendukung argumentasi penulis.

Contoh naskah ini menggunakan penyusunan dokumen sebagai konteks pembelajaran format. Penulis memisahkan isi naskah dari pengaturan tampilannya, sehingga perubahan judul bagian atau penambahan tabel dapat diikuti oleh pembaruan nomor dan rujukan. Pada penelitian sesungguhnya, konteks ini diganti dengan persoalan substantif yang menjadi fokus skripsi. Istilah teknis yang belum dikenali pembaca perlu dijelaskan sebelum digunakan untuk membangun argumen.

Penutup latar belakang sebaiknya menghubungkan masalah dengan rencana jawaban penelitian. Hubungan itu diperjelas melalui pertanyaan, tujuan, serta batasan yang saling sesuai. Bab ini menunjukkan bentuk hubungan tersebut; Bab \ref{bab:metodologi} memperlihatkan cara menata tahapan, sedangkan Bab \ref{bab:hasil} menunjukkan tempat penyajian keluaran. Rujukan silang tersebut akan mengikuti perubahan nomor bab secara otomatis.

\section{Rumusan Masalah}

Rumusan masalah ditulis sebagai pertanyaan yang dapat dijawab oleh prosedur dan bukti penelitian. Pada contoh penulisan ini, pertanyaannya digunakan untuk mengarahkan penggunaan template:

\begin{enumerate}
  \item Bagaimana menyusun naskah dengan hubungan yang jelas antara pendahuluan, landasan, metodologi, hasil, pembahasan, dan penutup?
  \item Bagaimana menyajikan tabel, gambar, persamaan, sitasi, dan lampiran yang mempunyai rujukan silang konsisten?
  \item Bagaimana memeriksa kesesuaian keluaran dokumen dengan sumber yang dapat disunting kembali?
\end{enumerate}

Pengguna perlu mengganti ketiga pertanyaan tersebut dengan rumusan masalah penelitiannya. Jumlah pertanyaan tidak perlu mengikuti jumlah pada contoh. Setiap pertanyaan harus memiliki kebutuhan bukti yang jelas; pertanyaan yang hanya menjelaskan kegiatan pengerjaan dapat dipertimbangkan kembali apabila tidak menghasilkan jawaban ilmiah.

\section{Tujuan}

Tujuan penulisan contoh ini adalah menyediakan susunan naskah yang dapat dipelajari dan disesuaikan. Keterkaitannya dengan rumusan masalah diperlihatkan pada Tabel \ref{tab:contoh-keterkaitan}. Tabel ini tidak menunjukkan capaian penelitian. Pengguna dapat memakai pola yang sama untuk memastikan setiap tujuan memiliki cara pemenuhan dan tempat pelaporan yang sesuai.

\begin{table}[htbp]
  \centering
  \caption{Contoh keterkaitan pertanyaan, tujuan, dan bagian pelaporan}
  \label{tab:contoh-keterkaitan}
  \begin{tabular}{p{3.6cm}p{4.2cm}p{4.4cm}}
    \toprule
    Fokus pertanyaan & Tujuan penulisan contoh & Bagian pelaporan \\
    \midrule
    Susunan naskah & Menunjukkan hubungan antarbagian & Pendahuluan sampai penutup \\
    Unsur dokumen & Menunjukkan rujukan silang & Tabel, gambar, persamaan, dan lampiran \\
    Pemeriksaan keluaran & Menjelaskan cara membaca keluaran & Pembahasan dan lampiran \\
    \bottomrule
  \end{tabular}
\end{table}

\section{Manfaat}

Manfaat contoh naskah bagi penulis adalah tersedianya titik awal yang mempunyai struktur dan unsur dokumen yang dapat diamati. Penulis dapat mengganti materi sambil mempertahankan pola penomoran serta rujukan yang diperlukan. Bagi pembaca, contoh ini memperlihatkan hubungan antara suatu unsur dokumen dan penjelasan yang menyertainya. Manfaat tersebut bersifat pembelajaran tata tulis, bukan klaim manfaat dari suatu metode penelitian.

Pada skripsi sesungguhnya, manfaat perlu dijelaskan menurut pihak yang menerima kegunaan hasil. Manfaat teoretis dapat dikaitkan dengan pengetahuan yang diperoleh, sedangkan manfaat praktis dapat dikaitkan dengan keputusan atau kegiatan yang terbantu. Pernyataan manfaat tidak sebaiknya menjanjikan sesuatu di luar ruang lingkup dan bukti penelitian.

\section{Batasan}

Contoh ini memperlihatkan struktur enam bab untuk susunan nonimplementatif. Susunan tersebut dapat disesuaikan dengan jenis penelitian, pedoman yang berlaku, dan arahan pembimbing. Nama bagian serta redaksi administratif perlu diperiksa kembali sebelum naskah digunakan sebagai dokumen akademik. Persetujuan atas isi, format akhir, dan kelengkapan administratif tidak timbul hanya karena dokumen dapat dikompilasi.

Isi contoh dibatasi pada penataan naskah. Tidak ada dataset, partisipan, skor evaluasi, ataupun hasil eksperimen yang disertakan. Tabel hasil menggunakan nama komponen dan deskripsi keluaran, sehingga tidak menyajikan angka empiris yang dibuat untuk mengisi halaman. Contoh rumus merupakan penjelasan notasi matematika dan dapat diganti dengan rumus yang relevan terhadap penelitian pengguna.

\section{Sistematika Penulisan}

Bab \ref{bab:pendahuluan} memperkenalkan konteks, pertanyaan, tujuan, manfaat, serta batas contoh. Bab \ref{bab:landasan} menyajikan konsep penataan dokumen dan contoh sitasi sumber resmi. Bab \ref{bab:metodologi} menunjukkan susunan prosedur dan diagram alur. Bab \ref{bab:hasil} memperlihatkan format pelaporan tanpa mengklaim hasil empiris. Bab \ref{bab:pembahasan} menunjukkan hubungan bukti dan penafsiran, sedangkan Bab \ref{bab:penutup} menata simpulan serta saran. Materi pendukung ditempatkan pada Lampiran \ref{lamp:contoh}.
